\section{The next integration and iteration}
\label{gs17-next}

For a logarithmic cutoff $R(g)=(\log(1/g))^p$ with $p>1/2$,
the Gaussian bad-block factor is smaller than any fixed power of
$g$, while $gR(g)^3$ tends to zero. These two estimates are
compatible: the same cutoff can suppress bad blocks and make
analytic cubic action remainders small. Jacobian terms that begin
linearly require their own local bounds; they are not cubic merely
because the action remainder is cubic.

The quantum step requires independent fixed-dimensional integration
coordinates, their normalized covariance, and localized estimates
for the complete interaction. The flat pivot covariance above
supplies such coordinates on the homogeneous regulator. The
singular value/gradient vector is useful for tests, but is not an
independent cell coordinate system for a Gaussian density expansion.
For complex continuation the tests used here remain fixed on the
real integration variables.

Notebook 18 gives a convergent cutoff expansion with overlapping
interaction supports under an explicit weighted coefficient bound.
Notebook 19 identifies the nonflat physical Gaussian and bounds
the local straightening, Haar and gauge factors. Notebook 20
verifies the complete coefficient and precision hypotheses for
the bare Wilson small-field step, with all normalization terms.
Notebook 21 supplies strict localization of the minimized
classical action uniformly in earlier depth and closes one
integration of that reference action. Later quantum steps
require transport to the next coarse background and control
of the accumulated interactions. Bad regions require
regional constructions and their normalization estimates. Iterating
those bounds must also control the running couplings and preserve
the observable hierarchy needed by the continuum construction.

The output sought from the iteration remains a common nontrivial
gauge-invariant continuum hierarchy with the Euclidean covariance,
positivity, regularity and Yang--Mills identification specified in
Notebook 01. The gap argument is a further target after that
existence construction is closed.

\begin{thebibliography}{9}
\bibitem{prior17}
\emph{Toward a four-dimensional Yang--Mills construction},
research Notebooks 5--16 in this workspace, 16 October 2026.
The flat reduced covariance, constrained Maxwell inverse,
nonlinear minimum and classical response theorems are used here.
\bibitem{source17}
E.~Sussman and collaborators,
\href{https://github.com/ethanwsussman/Yang--Mills-LLM-experiment}{Yang--Mills-LLM-experiment},
deposited raw accounts at commit
\texttt{e392ef45013a581664bab76b90867170b30fe10e}.
The Yang--Mills files are unchanged at checked commit
\texttt{785a68fdeef78a1da854c4807fe1f17b597d05e7}.
\bibitem{balaban17}
T.~Ba\l aban, \emph{Renormalization group approach to lattice gauge
field theories I}, Commun. Math. Phys. \textbf{109} (1987), 249--301,
\href{https://scispace.com/pdf/renormalization-group-approach-to-lattice-gauge-field-3ygai0vxt1.pdf}{original article copy}.
The source-coordinate match is discussed explicitly above.
\end{thebibliography}
\end{document}
